\subsection{Operational Semantics of \GRFRP}
\label{sec:semantics:grfrp:opsem}

\GRFRP services interact with the broader system \via an interface to the communication bus,
through which they import and export flows. To minimize messages on the bus, services employ
a push-based strategy, inspired by \frp~\cite{DBLP:conf/haskell/Elliott09}, propagating only signal
updates and event occurrences. When an event occurs or a signal is updated within a service, all
dependent flows are recomputed to maintain their values up-to-date, initiating a
deterministic cascade of computations. The operational semantics, described below,
captures this propagation process.

\begin{table}[!ht]
    \centering
    \begin{tabularx}{\textwidth}{|l|X|}\hline
        $v, \, \some{v}, \, \none \in \Val$                                                        & A value of any type; \ttf{some} and \none represent optional values.                                      \\\hline
        $v_\bot \in \ValBot ::= v \, \vert \, \bot$                                                & A value that may be present $v$ or absent $\bot$.                                                         \\\hline
        $v_d$                                                                                      & A default value, predefined unique value for each type.                                                   \\\hline
        $l ::= v::l \, \vert \, [v^+] \, \vert \, l_1 \lstconcat l_2$                              & An implicitly finite list of values $v^+$.                                                                \\\hline
        $sty ::= ty \, \vert \, \mayB{ty}$                                                         & A type, where \kwf{?} indicates an event type.                                                            \\\hline
        $\tyCtx[x] = sty$                                                                          & A typing context \tyCtx mapping flow identifiers to their types.                                          \\\hline
        \mkcell{$\G[f] = \mathtiny{\component{f}{\mathtiny{x_i^+}}{\mathtiny{x_o^+}}{\mathtiny{eq^+}}{\mathtiny{x_m^+ = c^+}}}$                                                                                \\
        $\G[F] = \service{F}{stmt^+}$}                                                             & A program context \G mapping component and service names to their definitions.                            \\\hline
        $i, \, i_T$                                                                                & Identifiers for flows and timers, respectively.                                                           \\\hline
        $\mess ::= \inputMess{i}{v} \, \vert \, \timerMess{i_T}$                                   & A message from the environment, either a flow input or a timer signal.                                    \\\hline
        $s$                                                                                        & The internal state of a service, a statement, or a flow operation.                                        \\\hline
        $\tau ::= (i_T, t) \diamond \tau \, \vert \, \queue{(i_T,t)^+}$                            & A timer queue consisting of pairs of timer identifiers $i_T$ and timestamps $t$.                          \\\hline
        \mkcell{$\gamma[x] = v_\bot$                                                                                                                                                                           \\
        $\gamma ::= [(v_\bot/x)^+] \, \vert \, \gamma_1 \ctxconcat \gamma_2$}                      & A propagation context $\gamma$ mapping flow identifiers to their updates or occurrences.                  \\\hline
        \stateTy{\epsilon}                                                                         & The type of compatible states for the \GRFRP construct $\epsilon$.                                        \\\hline
        $\initstate{\epsilon} \in \stateTy{\epsilon}$                                              & The initial state of the \GRFRP construct $\epsilon$.                                                     \\\hline
        \timerTy{\epsilon}                                                                         & The type of compatible list (queue for services) of timers for the \GRFRP construct $\epsilon$.           \\\hline
        $\inittimer{\epsilon} \in \timerTy{\epsilon}$                                              & The initial list (queue for services) of timers of the \GRFRP construct $\epsilon$.                       \\\hline
        \mathfoot{\serviceOpSem{s}{\tau}{F}{\mess}{t}{s'}{\tau'}{(x_o, v_o)^+}}                    & The service $F$ reacts to a message at time $t$, updating its state and producing outputs.                \\\hline
        \mathtiny{\stmtOpSem{\gamma}{s}{stmt^+}{i, t, v}{\gamma'}{s'}{(x_o, v_o)^+}{(i_T, t_T)^+}} & Propagation of an update through flow statements, producing new states, output updates, and timer resets. \\\hline
        \mathfoot{\sigmaOpSem{\gamma}{s}{\sigma}{t}{s'}{v_\bot^+}{(i_T, t_T)^+}}                   & Evaluation of $\sigma$ at time $t$, producing a new state, outputs, and timer resets.                     \\\hline
    \end{tabularx}
    \caption[Notations for the operational semantics of \GRFRP.]%
    {Notations for the operational semantics of \GRFRP.}
    \label{tab:semantics:grfrp:opsem:notations}
\end{table}

We first introduce the notations used throughout this semantics, summarized in
\Cref{tab:semantics:grfrp:opsem:notations}.
%
The set of all possible values is denoted by \Val, a value is denoted by $v$, and may be wrapped as
an optional using \some{v} or \none.
%
A bottom value $v_\bot$ can be either a defined value $v$ or the undefined marker $\bot$, used to
represent the presence or absence of a value. The set of all bottom values is then the addition of
\Val with the singleton $\{\bot\}$, denoted by \ValBot.
%
As signals always have a value, they start with a default value $v_d$, which is a predefined unique
value for each type (\eg \ttf{0} for integers and \ttf{0.} for floats). This default value
represents the value of a signal before any explicit updates occur.
%
Across the entire semantics, we use the notation $X^+$ to denote a finite sequence of $X$, this sequence can be empty.
Lists of elements $v^+$ are denoted by $[v^+]$, may be inductively constructed as $v::l$, and
concatenated as $l_1 \lstconcat l_2$.
%
The sampling types $sty$ of flows (\ie the type of the value visible in a flow at any given time)
are recorded in a typing context \tyCtx. A signal type is denoted by $ty$, while an event type is
denoted by \mayB{ty}, indicating that an event may have no value. The sampling type of a flow
identified by $x$ is accessed as $\tyCtx[x] = sty$. These types serve as syntactic labels, but their
correct usage is governed by a type system, detailed in \Cref{sec:grust:typing} of \Cref{cha:grust}.
In the semantics, they are used to distinguish the behavior of signals from that of events.
%
A program is represented as a mapping \G from component and service names to their definitions.
The definition of a component $f$ is accessed by $\G[F]$, and the one of a service $F$ by $\G[F]$.
%
Identifiers for imported flows and timers are written $i$ and $i_T$, respectively.
%
Messages received by a service from the environment are denoted by \mess and are of the form
\inputMess{i}{v} for a new input value $v$ on flow $i$ (not necessarily imported by the service), or
\timerMess{i_T} for a timer expiration on timer $i_T$.
%
States for services, statements, flow operations are written $s$.
%
Timer queues $\tau$ represent scheduled timers, and may be expressed either as a sequence
\queue{(i_T,t)^+} or inductively using $(i_T, t) \diamond \tau$, indicating that the timer
$i_T$ is scheduled to trigger at time $t$, followed by the remaining queue $\tau$.
%
The propagation context $\gamma$ tracks whether each flow has been updated or triggered. For a
signal $x_s$, $\gamma[x_s] = \bot$ indicates that the signal has not changed. For an event $x_e$,
$\gamma[x_e] = \bot$ denotes the absence of an occurrence at the current step. Otherwise,
$\gamma[x] = v$ represents the newly updated signal value or the emitted event payload.
%
The concatenation of two update contexts $\gamma_1$ and $\gamma_2$, is defined as their disjoint
union, $\gamma_1 \ctxconcat \gamma_2$. An update context mapping bottom values $v_\bot^+$ to
distinct identifiers $x^+$ is denoted by $[(v_\bot/x)^+]$.
%
The constructs of \GRFRP are designed to react to event occurrences, signal updates, and timer
expirations. Each reaction triggers a computation step that can modify the internal state (\eg, by
retaining the last value of a signal) and reset timers.
%
For all \GRFRP construct $\epsilon$, the set of all states corresponding to $\epsilon$ is denoted by
\stateTy{\epsilon} and its initial state is denoted by \initstate{\epsilon}; the set of all lists
(queues for services) of times corresponding to $\epsilon$ is denoted by \timerTy{\epsilon} and its
initial timers are denoted by \inittimer{\epsilon}. The following paragraphs explain the computation
steps of the different constructs.

\paragraph{Receive message}

Consider a service $F$ defined in a program \G within a typing context \tyCtx, and let $s$ and
$\tau$ respectively denote the state and timer queue of $F$. The judgment
\begin{equation*}
    \serviceOpSem{s}{\tau}{F}{\mess}{t}{s'}{\tau'}{(x_o, v_o)^+}
\end{equation*}
describes the reaction of service $F$ to a message \mess received at time $t$. This transition
results in a new state $s'$, an updated timer queue $\tau'$, and a list of output updates
$[(x_o, v_o)^+]$ that are transmitted on the communication bus.

A state corresponding to the service $F$ is a state corresponding to its statements
(\Cref{eq:semantics:grfrp:opsem:init:service:state}).
The initial state \initstate{F} and initial timer queue
\inittimer{F} are computed from the initializations of each statement, as defined
in \Crefrange{eq:semantics:grfrp:opsem:init:service:state}{eq:semantics:grfrp:opsem:init:service:timer},
which use the function \smf{reset} defined by \Cref{rule:semantics:grfrp:opsem:aux:reset}.
\begin{align}
    \stateTy{F}    & = \stateTy{stmt^+}                                        & \initstate{F} & = \initstate{stmt^+} \label{eq:semantics:grfrp:opsem:init:service:state}                       \\
    \timerTy{F}    & = \Queue{U \times \Time}                                  & \inittimer{F} & = \reset{\emptytimer}{\inittimer{stmt^+}}  \label{eq:semantics:grfrp:opsem:init:service:timer} \\
    \text{with} \; & \timerTy{stmt^+} = \List{U \times \Time} \; \text{and} \; & \G[F]         & = \service{F}{stmt^+} \nonumber
\end{align}

Message handling follows two semantic rules, corresponding to the two kinds of messages a service
may receive:
\nameref{rule:semantics:grfrp:opsem:step:service:input} for input messages \inputMess{i}{v} from
the communication bus;
and \nameref{rule:semantics:grfrp:opsem:step:service:timer} for internal timer expiration messages
\timerMess{i_T}.
%
In both cases, the message is propagated through the statements of the service ($stmt^+$) using the
corresponding identifier ($i$ or $i_T$), timestamp $t$, and value ($v$ or $()$). This propagation,
described in the following paragraph, produces:
an updated context $\gamma$;
a new internal state $s'$;
a list of output updates $[(x_o, v_o)^+]$;
and a list of timers to be reset $[(i_T, t_T)^+]$.
%
The timer queue $\tau$ is updated using the function \smf{reset}, defined in
\Cref{rule:semantics:grfrp:opsem:aux:reset}, which replaces or inserts the given timers in the
queue so that the queue remains increasing in timestamp.

\begin{center}
    \begin{minipage}{0.9\textwidth}
        \centering
        \begin{mathpar}
            %%%%% Runtime handles input
            \inferOPSEMMess
            {
                \G[F] = \service{F}{stmt^+} \\
                (\tau = (i_T, t_T) \diamond \tau' \wedge t_i \le t_T) \vee (\tau = \emptytimer) \\
                \stmtOpSem{\emptygamma}{s}{stmt^+}{i, t_i, v_i}{\gamma}{s'}{(x_o, v_o)^+}{(i_T, t_T)^+}
            }
            {
                \serviceOpSem{s}{\tau}{F}{\inputMess{i}{v_i}}{t_i}{s'}{\reset{\tau}{[(i_T, t_T)^+]}}{(x_o, v_o)^+}
            }
            {Input}
            \label{rule:semantics:grfrp:opsem:step:service:input}

            %%%%% Runtime handles timer
            \inferOPSEMMess
            {
                \G[F] = \service{F}{stmt^+} \\
                \tau = (i_T, t_T) \diamond \tau' \\
                \stmtOpSem{[()/i_T]}{s}{stmt^+}{i_T, t_T, ()}{\gamma}{s'}{(x_o, v_o)^+}{(i_T, t_T)^+}
            }
            {
                \serviceOpSem{s}{\tau}{F}{\timerMess{i_T}}{t_T}{s'}{\reset{\tau'}{[(i_T, t_T)^+]}}{(x_o, v_o)^+}
            }
            {Timer}
            \label{rule:semantics:grfrp:opsem:step:service:timer}
        \end{mathpar}
    \end{minipage}
\end{center}

\begin{remark}
    Rule \nameref{rule:semantics:grfrp:opsem:step:service:input} initiates the propagation through
    the service's statements with an empty context (\emptygamma). This reflects the fact that the
    communication bus transmits all flows in the system, including those not imported or used by the
    service. The operational semantics of the import statement (described in the following
    paragraph) is responsible for selecting the relevant input flows, adding their values to the
    update context, and propagating them through the service.

    In contrast, timers are managed internally by the service using the timer queues $\tau$. Timer
    occurrences are always relevant to the service and are therefore directly added to the context
    and propagated through the service's statements.

    Together, rules \nameref{rule:semantics:grfrp:opsem:step:service:input} and
    \nameref{rule:semantics:grfrp:opsem:step:service:timer} ensure two key properties: (1) timers
    are triggered only after all inputs from the communication bus have been processed, and (2)
    timestamps are handled in ascending order. While this ascending order of timestamps is
    fundamental to the semantics, access to a global time and its use as a computation timestamp
    could enable an alternative message prioritization scheme (\eg based on the ASIL level of the
    message).
\end{remark}

\begin{align}
    \reset{\tau}{l}           & = \ttf{fold}(\lambda (\tau_\textit{acc}, \, (i_T, t_T)). \resetAux{\tau_\textit{acc}}{i_T}{t_T}, \, \tau, \, l)  \label{rule:semantics:grfrp:opsem:aux:reset} \\
    \resetAux{\tau}{i_T}{t_T} & =  \ttf{match} \, \tau \, \{ \label{rule:semantics:grfrp:opsem:aux:reset:aux}                                                                                 \\
                              & \quad \qquad \begin{array}{lll}
                                                 \emptytimer                    &                           & \rightarrow \queue{(i_T, t_T)}                                  \\
                                                 (i_{T'}, t_{T'}) \diamond \tau & \text{if} \, i_{T'} = i_T & \rightarrow \resetAux{\tau}{i_T}{t_T}                           \\
                                                 (i_{T'}, t_{T'}) \diamond \tau & \text{if} \, t_{T'} < t_T & \rightarrow (i_{T'}, t_{T'}) \diamond \resetAux{\tau}{i_T}{t_T} \\
                                                 (i_{T'}, t_{T'}) \diamond \tau &                           & \rightarrow (i_T, t_T) \diamond (i_{T'}, t_{T'}) \diamond \tau  \\
                                             \end{array} \nonumber                     \\
                              & \; \quad \} \nonumber
\end{align}

For example, consider the \idf{aeb} service from \Cref{fig:semantics:example:aeb} and the execution
depicted in the chronogram of \Cref{fig:semantics:example:aeb:chronogram}. We focus on the
transition step of the service at time $t + 1s$, when it receives the event \idf{pedestrian\_l}.
%
Prior to the transition, the service is in the state \No with a timer queue
\queue{(\idf{t_p}, t_0)}, where \idf{t_p} is the timer associated with the \kwf{timeout} operator
from \Cref{ex:line:timeout}, and $t_0 > t + 1s$.
%
Upon receiving the input message \inputMess{\idf{pedestrian\_l}}{20.} at time $t + 1s$, the service
transitions to a new timer queue \queue{(\idf{t_p}, t+3s)}, emits the strategy \Soft on the output
\idf{strat}, and updates its internal state to \Soft.
%
This transition step in the semantics is expressed as:
\begin{equation*}
    (\No, \queue{(\idf{t_p},t_0)})
    \xRightarrow[\service{\kwf{aeb}}{stmt^+}]{\inputMess{\idf{pedestrian\_l}}{20.}, \, t+1s}
    (\Soft, \queue{(\idf{t_p},t+3s)}): [(\idf{strat}, \Soft)].
\end{equation*}
%
This transition follows the rule \nameref{rule:semantics:grfrp:opsem:step:service:input}, which
requires a transition of the service's statements as follows:
\begin{equation*}
    (\emptygamma, \No)
    \xRightarrow[stmt^+]{\idf{pedestrian\_l}, \, t+1s, \, 20.}
    (\mathtiny{\left[\begin{array}{l}
                20./\idf{pedestrian\_l},  \\
                20./\idf{pedestrian},     \\
                \bot/\idf{t\_pedestrian}, \\
                \Soft/\idf{strat}
            \end{array}\right]}, \Soft):
    ([(\idf{strat}, \Soft)], \, [(\idf{t_p}, t+3s)]).
\end{equation*}
%
The intermediate steps leading to this transition will be detailed in the following paragraphs.

\paragraph{Propagate information in statements}

Consider a sequence of statements $stmt^+$ defined in a program \G and typing context \tyCtx, and
let $\gamma$ and $s$ respectively denote the current value context of the propagation and the state
of $stmt^+$. The judgment
\begin{equation*}
    \stmtOpSem{\gamma}{s}{stmt^+}{i, t, v}{\gamma'}{s'}{(x_o, v_o)^+}{(i_T, t_T)^+}
\end{equation*}
describes the propagation of a new value $v$ for identifier $i$ at time $t$ through
$stmt^+$ starting with the context $\gamma$. This propagation results in an updated context
$\gamma'$, a new state $s'$, a list of output updates $[(x_o, v_o)^+]$, and a list of timer resets
$[(i_T, t_T)^+]$.

\subparagraph{Imports}
An import statement contributes nothing to the overall state of the block of statements, nor to the
initial timer queue, as defined by \Crefrange{eq:semantics:grfrp:opsem:stateTy:stmt:import}
{eq:semantics:grfrp:opsem:init:stmt:import:timer}.
\begin{align}
    \stateTy{\import{x}, stmt^+}   & = \stateTy{stmt^+} \label{eq:semantics:grfrp:opsem:stateTy:stmt:import}      \\
    \initstate{\import{x}, stmt^+} & = \initstate{stmt^+} \label{eq:semantics:grfrp:opsem:init:stmt:import:state} \\
    \timerTy{\import{x}, stmt^+}   & = \timerTy{stmt^+} \label{eq:semantics:grfrp:opsem:timerTy:stmt:import}      \\
    \inittimer{\import{x}, stmt^+} & = \inittimer{stmt^+} \label{eq:semantics:grfrp:opsem:init:stmt:import:timer}
\end{align}

The rule \nameref{rule:semantics:grfrp:opsem:step:stmt:import:present} handles the import of the
flow $i$ that triggered the propagation. Its value $v$ is added to the update context $\gamma$, and
the propagation continues.\footnote{Once the input message has been imported, its value no longer
    needs to be propagated; only the update context $\gamma \ctxconcat [v/x]$ is required. However,
    the value is retained for convenience in writing.}
%
If the imported identifier $x$ does not correspond to the propagated flow, rule
\nameref{rule:semantics:grfrp:opsem:step:stmt:import:absent} applies, passing the update context
$\gamma$ unchanged to the next statements.

\begin{center}
    \begin{minipage}{0.9\textwidth}
        \centering
        \begin{mathpar}
            % ImportPresent
            \inferOPSEMFlowStmt
            {
                x = i \\
                \stmtOpSem{\gamma \ctxconcat [v/x]}{s}{stmt^+}{i, t, v}{\gamma'}{s'}{(x_o, v_o)^+}{(i_T, t_T)^+}
            }
            {
                \stmtOpSem{\gamma}{s}{\import{x}, stmt^+}{i, t, v}{\gamma'}{s'}{(x_o, v_o)^+}{(i_T, t_T)^+}
            }
            {Import{\scriptsize pres}}
            \label{rule:semantics:grfrp:opsem:step:stmt:import:present}

            % ImportAbsent
            \inferOPSEMFlowStmt
            {
                x \neq i \\
                \stmtOpSem{\gamma}{s}{stmt^+}{i, t, v}{\gamma'}{s'}{(x_o, v_o)^+}{(i_T, t_T)^+}
            }
            {
                \stmtOpSem{\gamma}{s}{\import{x}, stmt^+}{i, t, v}{\gamma'}{s'}{(x_o, v_o)^+}{(i_T, t_T)^+}
            }
            {Import{\scriptsize abs}}
            \label{rule:semantics:grfrp:opsem:step:stmt:import:absent}
        \end{mathpar}
    \end{minipage}
\end{center}

Continuing with the \aeb example from \Cref{fig:semantics:example:aeb}, we again focus on the
transition step of the \idf{aeb} service at time $t + 1s$, as depicted in the chronogram of
\Cref{fig:semantics:example:aeb:chronogram}. Here, we detail the handling of the three import
statements:
\begin{multline*}
    (\emptygamma, \No)
    \xRightarrow[\mathtiny{
            \import{\idf{speed}},
            \import{\idf{pedestrian\_l}},
            \import{\idf{pedestrian\_r}},
            stmt^+}
    ]{\idf{pedestrian\_l}, \, t+1s, \, 20.}
    (\mathtiny{\left[\begin{array}{l}
                20./\idf{pedestrian\_l},  \\
                20./\idf{pedestrian},     \\
                \bot/\idf{t\_pedestrian}, \\
                \Soft/\idf{strat}
            \end{array}\right]}, \Soft): \\
    ([(\idf{strat}, \Soft)], \, [(\idf{t_p}, t+3s)]).
\end{multline*}
%
The first import, \idf{speed}, does not correspond to the identifier of the flow that triggered the
computation. It follows rule \nameref{rule:semantics:grfrp:opsem:step:stmt:import:absent}, which
requires a transition of the subsequent statements as follows:
\begin{multline*}
    (\emptygamma, \No)
    \xRightarrow[\mathtiny{
            \import{\idf{pedestrian\_l}},
            \import{\idf{pedestrian\_r}},
            stmt^+
        }]{\idf{pedestrian\_l}, \, t+1s, \, 20.}
    (\mathtiny{\left[\begin{array}{l}
                20./\idf{pedestrian\_l},  \\
                20./\idf{pedestrian},     \\
                \bot/\idf{t\_pedestrian}, \\
                \Soft/\idf{strat}
            \end{array}\right]}, \Soft): \\
    ([(\idf{strat}, \Soft)], \, [(\idf{t_p}, t+3s)]).
\end{multline*}
%
Next, the import of \idf{pedestrian\_l}, which is the identifier of the flow that triggered the
computation, follows rule \nameref{rule:semantics:grfrp:opsem:step:stmt:import:present}. This rule
requires a transition of the subsequent statements as follows:
\begin{multline*}
    ([\mathtiny{20./\idf{pedestrian\_l}}], \No)
    \xRightarrow[\mathtiny{
            \import{\idf{pedestrian\_r}},
            stmt^+
        }]{\idf{pedestrian\_l}, \, t+1s, \, 20.}
    (\mathtiny{\left[\begin{array}{l}
                20./\idf{pedestrian\_l},  \\
                20./\idf{pedestrian},     \\
                \bot/\idf{t\_pedestrian}, \\
                \Soft/\idf{strat}
            \end{array}\right]}, \Soft): \\
    ([(\idf{strat}, \Soft)], \, [(\idf{t_p}, t+3s)]).
\end{multline*}
%
Finally, the import of \idf{pedestrian\_r}, which also does not correspond to the identifier of the
flow that triggered the computation, is similar to the import of \idf{speed}. It follows rule
\nameref{rule:semantics:grfrp:opsem:step:stmt:import:absent} and requires a transition of the
subsequent statements as follows:
\begin{multline*}
    ([\mathtiny{20./\idf{pedestrian\_l}}], \No)
    \xRightarrow[stmt^+]{\idf{pedestrian\_l}, \, t+1s, \, 20.}
    (\mathtiny{\left[\begin{array}{l}
                20./\idf{pedestrian\_l},  \\
                20./\idf{pedestrian},     \\
                \bot/\idf{t\_pedestrian}, \\
                \Soft/\idf{strat}
            \end{array}\right]}, \Soft): \\
    ([(\idf{strat}, \Soft)], \, [(\idf{t_p}, t+3s)]).
\end{multline*}
%
The remaining steps leading to this transition will be described in the following paragraphs.

\subparagraph{Flow declarations}
The state corresponding to a block of statements $stmt^+$ preceded by a flow declaration
\decl{x^+}{\sigma}, as defined by \Cref{eq:semantics:grfrp:opsem:stateTy:stmt:decl}, is the
Cartesian product of the state corresponding to the flow operation $\sigma$ and the state
corresponding to $stmt^+$.
%
The initial state and timer queue of flow declarations include the initial state and timers of their
flow operations, as defined by \Crefrange{eq:semantics:grfrp:opsem:init:stmt:decl:state}
{eq:semantics:grfrp:opsem:init:stmt:decl:timer}.\footnote{The Cartesian product for states and
    concatenation for timers reflect the structure of the implementation: each flow operation state is
    represented as a distinct field in the global state; initial timers are added successively to the
    timer queue, and since each timer has a unique identifier, they do not interfere with each other.
    This ensures that each part of the system maintains its own state and timers independently.}
\begin{align}
    \stateTy{\decl{x^+}{\sigma}, stmt^+}   & = \stateTy{\sigma} \times \stateTy{stmt^+} \label{eq:semantics:grfrp:opsem:stateTy:stmt:decl}                    \\
    \initstate{\decl{x^+}{\sigma}, stmt^+} & = (\initstate{\sigma}, \initstate{stmt^+}) \label{eq:semantics:grfrp:opsem:init:stmt:decl:state}                 \\
    \timerTy{\decl{x^+}{\sigma}, stmt^+}   & = \List{U_1 \cup U_2 \times \Time} \label{eq:semantics:grfrp:opsem:timerTy:stmt:decl}                            \\
    \text{with} \;                         & \timerTy{\sigma} = \List{U_1 \times \Time} \; \text{and} \; \timerTy{stmt^+} = \List{U_2 \times \Time} \nonumber \\
    \inittimer{\decl{x^+}{\sigma}, stmt^+} & = \inittimer{\sigma} \lstconcat \inittimer{stmt^+} \label{eq:semantics:grfrp:opsem:init:stmt:decl:timer}
\end{align}

The operational semantics for flow declarations are governed by rule
\nameref{rule:semantics:grfrp:opsem:step:stmt:decl}. This rule first evaluates the declared flow
operation $\sigma$ at time $t$ under state $s_\sigma$, yielding a new state $s_\sigma'$, output
values $v_\bot^+$, and timer resets $(i_\sigma, t_\sigma)^+$. The context $\gamma$ is then extended
with the output values $(v_\bot/x)^+$ to form a new context for the subsequent statements. The timer
resets produced by the flow operation are combined with any timer resets produced by the subsequent
statements.

\begin{center}
    \begin{minipage}{0.95\textwidth}
        \centering
        \begin{mathpar}
            % Let
            \inferOPSEMFlowStmt
            {
                \sigmaOpSem{\gamma}{s_\sigma}{\sigma}{t}{s'_\sigma}{v_\bot^+}{(i_\sigma, t_\sigma)^+} \\
                \gamma_\sigma = \gamma \ctxconcat [(v_\bot/x)^+] \\
                \stmtOpSem{\gamma_\sigma}{s}{stmt^+}{i, t, v}{\gamma'}{s'}{(x_o, v_o)^+}{(i_T, t_T)^+}
            }
            {
                \stmtOpSem{\gamma}{(s_\sigma, s)}{\decl{x^+}{\sigma}, stmt^+}{i, t, v}{\gamma'}{(s'_\sigma, s')}{(x_o, v_o)^+}{(i_\sigma, t_\sigma)^+, (i_T, t_T)^+}
            }
            {Let}
            \label{rule:semantics:grfrp:opsem:step:stmt:decl}
        \end{mathpar}
    \end{minipage}
\end{center}

To illustrate this process, we continue with the transition step of the \idf{aeb} service at time
$t + 1s$. We now detail the three declaration statements, all of which follow rule
\nameref{rule:semantics:grfrp:opsem:step:stmt:decl}:
\begin{multline*}
    ([\mathtiny{20./\idf{pedestrian\_l}}], \No)
    \xRightarrow[\mathtiny{
            \decl{\idf{pedestrian}}{\mergeop{\idf{pedestrian\_l}}{\idf{pedestrian\_r}}},
            \decl{\idf{t\_pedestrian}}{\dots},
            \decl{\idf{strat}}{\dots},
            stmt^+
        }]{\idf{pedestrian\_l}, \, t+1s, \, 20.} \\
    (\mathtiny{\left[\begin{array}{l}
                20./\idf{pedestrian\_l},  \\
                20./\idf{pedestrian},     \\
                \bot/\idf{t\_pedestrian}, \\
                \Soft/\idf{strat}
            \end{array}\right]}, \Soft):
    ([(\idf{strat}, \Soft)], \, [(\idf{t_p}, t+3s)]).
\end{multline*}
%
The transitions of the flow operations themselves are detailed in the next paragraph; this example
focuses on their outcome.
%
First, the declaration of \idf{pedestrian} merges the occurrences of \idf{pedestrian\_l} and
\idf{pedestrian\_r} in the update context. Since \idf{pedestrian\_l} is present, the identifier
\idf{pedestrian} is mapped to its value in the context for the transition of the subsequent
statements:
\begin{multline*}
    (\mathtiny{\left[\begin{array}{l}
                20./\idf{pedestrian\_l}, \\
                20./\idf{pedestrian}
            \end{array}\right]}, \No)
    \xRightarrow[\mathtiny{
            \decl{\idf{t\_pedestrian}}{\timeout{\idf{pedestrian}}{2000}},
            \decl{\idf{strat}}{\dots},
            stmt^+
        }]{\idf{pedestrian\_l}, \, t+1s, \, 20.} \\
    (\mathtiny{\left[\begin{array}{l}
                20./\idf{pedestrian\_l},  \\
                20./\idf{pedestrian},     \\
                \bot/\idf{t\_pedestrian}, \\
                \Soft/\idf{strat}
            \end{array}\right]}, \Soft):
    ([(\idf{strat}, \Soft)], \, [(\idf{t_p}, t+3s)]).
\end{multline*}
%
Next, the declaration of \idf{t\_pedestrian}, which is the timeout of the event \idf{pedestrian},
triggers a reset of the timer \idf{t_p} to $t+3s$. According to rule
\nameref{rule:semantics:grfrp:opsem:step:stmt:decl}, this timer reset is produced solely by the flow
operation within this declaration and is then combined with any timer resets produced by the next
statements (here, there are none):
\begin{multline*}
    (\mathtiny{\left[\begin{array}{l}
                20./\idf{pedestrian\_l}, \\
                20./\idf{pedestrian},    \\
                \bot/\idf{t\_pedestrian}
            \end{array}\right]}, \No)
    \xRightarrow[\mathtiny{
            \decl{\idf{strat}}{\idf{braking\_strat}(\idf{pedestrian}, \idf{t\_pedestrian}, \idf{speed})},
            stmt^+
        }]{\idf{pedestrian\_l}, \, t+1s, \, 20.} \\
    (\mathtiny{\left[\begin{array}{l}
                20./\idf{pedestrian\_l},  \\
                20./\idf{pedestrian},     \\
                \bot/\idf{t\_pedestrian}, \\
                \Soft/\idf{strat}
            \end{array}\right]}, \Soft):
    ([(\idf{strat}, \Soft)], \, []).
\end{multline*}
%
Finally, the declaration of \idf{strat} results from the application of the component
\idf{braking\_strat} with the inputs \idf{pedestrian}, \idf{t\_pedestrian}, and \idf{speed}, leading
to the \Soft braking strategy. The state value transitions from \No to \Soft during this component
application, but this state change is local to this declaration and is not reflected in the
subsequent statements:
\begin{equation*}
    (\mathtiny{\left[\begin{array}{l}
                20./\idf{pedestrian\_l},  \\
                20./\idf{pedestrian},     \\
                \bot/\idf{t\_pedestrian}, \\
                \Soft/\idf{strat}
            \end{array}\right]}, ())
    \xRightarrow[stmt^+]{\idf{pedestrian\_l}, \, t+1s, \, 20.}
    (\mathtiny{\left[\begin{array}{l}
                20./\idf{pedestrian\_l},  \\
                20./\idf{pedestrian},     \\
                \bot/\idf{t\_pedestrian}, \\
                \Soft/\idf{strat}
            \end{array}\right]}, ()):
    ([(\idf{strat}, \Soft)], \, []).
\end{equation*}
%
The remaining steps leading to this transition will be described in the following paragraphs.

\subparagraph{Exports}
An export statement contributes nothing to the overall state of the block of statements, nor to the
initial timer queue, as defined by \Crefrange{eq:semantics:grfrp:opsem:stateTy:stmt:export}
{eq:semantics:grfrp:opsem:init:stmt:export:timer}.
\begin{align}
    \stateTy{\export{x}, stmt^+}   & = \stateTy{stmt^+} \label{eq:semantics:grfrp:opsem:stateTy:stmt:export}      \\
    \initstate{\export{x}, stmt^+} & = \initstate{stmt^+} \label{eq:semantics:grfrp:opsem:init:stmt:export:state} \\
    \timerTy{\export{x}, stmt^+}   & = \timerTy{stmt^+} \label{eq:semantics:grfrp:opsem:timerTy:stmt:export}      \\
    \inittimer{\export{x}, stmt^+} & = \inittimer{stmt^+} \label{eq:semantics:grfrp:opsem:init:stmt:export:timer}
\end{align}

If the value $\gamma[x] = v$ is present, rule
\nameref{rule:semantics:grfrp:opsem:step:stmt:export:value} appends $(x, v)$ to the output list.
If $\gamma[x] = \bot$, rule \nameref{rule:semantics:grfrp:opsem:step:stmt:export:bottom} simply
skips the export and proceeds. Therefore, only flow changes/occurrences are exported to the
communication bus.

\begin{center}
    \begin{minipage}{0.9\textwidth}
        \centering
        \begin{mathpar}
            % Exportval
            \inferOPSEMFlowStmt
            {
                \gamma[x] = v \\
                \stmtOpSem{\gamma}{s}{stmt^+}{i, t, v}{\gamma'}{s'}{(x_o, v_o)^+}{(i_T, t_T)^+}
            }
            {
                \stmtOpSem{\gamma}{s}{\export{x}, stmt^+}{i, t, v}{\gamma'}{s'}{(x_o, v_o)^+, (x, v)}{(i_T, t_T)^+}
            }
            {Export{\scriptsize val}}
            \label{rule:semantics:grfrp:opsem:step:stmt:export:value}

            % ExportBottom
            \inferOPSEMFlowStmt
            {
                \gamma[x] = \bot \\
                \stmtOpSem{\gamma}{s}{stmt^+}{i, t, v}{\gamma'}{s'}{(x_o, v_o)^+}{(i_T, t_T)^+}
            }
            {
                \stmtOpSem{\gamma}{s}{\export{x}, stmt^+}{i, t, v}{\gamma'}{s'}{(x_o, v_o)^+}{(i_T, t_T)^+}
            }
            {Export{\scriptsize bot}}
            \label{rule:semantics:grfrp:opsem:step:stmt:export:bottom}
        \end{mathpar}
    \end{minipage}
\end{center}

We continue with the transition step of the \idf{aeb} service at time $t + 1s$, detailing the
transition of the export statement:
\begin{equation*}
    (\mathtiny{\left[\begin{array}{l}
                20./\idf{pedestrian\_l},  \\
                20./\idf{pedestrian},     \\
                \bot/\idf{t\_pedestrian}, \\
                \Soft/\idf{strat}
            \end{array}\right]}, ())
    \xRightarrow[\export{\idf{start}}]{\idf{pedestrian\_l}, \, t+1s, \, 20.}
    (\mathtiny{\left[\begin{array}{l}
                20./\idf{pedestrian\_l},  \\
                20./\idf{pedestrian},     \\
                \bot/\idf{t\_pedestrian}, \\
                \Soft/\idf{strat}
            \end{array}\right]}, ()):
    ([(\idf{strat}, \Soft)], \, []).
\end{equation*}
%
Because \idf{strat} is present in the update context, it follows rule
\nameref{rule:semantics:grfrp:opsem:step:stmt:export:value}, which is responsible for the output
\idf{strat}. This output is then removed from the transition of the subsequent statements
(the skip statement \skp in this case):
\begin{equation*}
    (\mathtiny{\left[\begin{array}{l}
                20./\idf{pedestrian\_l},  \\
                20./\idf{pedestrian},     \\
                \bot/\idf{t\_pedestrian}, \\
                \Soft/\idf{strat}
            \end{array}\right]}, ())
    \xRightarrow[\skp]{\idf{pedestrian\_l}, \, t+1s, \, 20.}
    (\mathtiny{\left[\begin{array}{l}
                20./\idf{pedestrian\_l},  \\
                20./\idf{pedestrian},     \\
                \bot/\idf{t\_pedestrian}, \\
                \Soft/\idf{strat}
            \end{array}\right]}, ()):
    ([], \, []).
\end{equation*}
%
The remaining step is described in the following paragraph.

\subparagraph{Skip}
Rule \nameref{rule:semantics:grfrp:opsem:step:stmt:skip} and
\Cref{eq:semantics:grfrp:opsem:init:stmt:skip} capture the base case where no statements remain. It
leaves an empty state, preserves the current context, and returns no outputs nor timer resets.
%
\begin{align}
    \stateTy{\skp}   & = \{()\}                                                     &
    \initstate{\skp} & = ()                                                         &
    \timerTy{\skp}   & = \List{\emptyset \times \Time}                              &
    \inittimer{\skp} & = \emptylist \label{eq:semantics:grfrp:opsem:init:stmt:skip}
\end{align}

\begin{center}
    \begin{minipage}{0.9\textwidth}
        \centering
        \begin{mathpar}
            % Skip
            \inferOPSEMFlowStmt
            { }
            { \stmtOpSem{\gamma}{()}{\skp}{i, t, v}{\gamma}{()}{\,}{\,} }
            {Skip}
            \label{rule:semantics:grfrp:opsem:step:stmt:skip}
        \end{mathpar}
    \end{minipage}
\end{center}

For instance, the rule \nameref{rule:semantics:grfrp:opsem:step:stmt:skip} finalizes the transition
step of the \idf{aeb} service at time $t + 1s$:
\begin{equation*}
    (\mathtiny{\left[\begin{array}{l}
                20./\idf{pedestrian\_l},  \\
                20./\idf{pedestrian},     \\
                \bot/\idf{t\_pedestrian}, \\
                \Soft/\idf{strat}
            \end{array}\right]}, ())
    \xRightarrow[\skp]{\idf{pedestrian\_l}, \, t+1s, \, 20.}
    (\mathtiny{\left[\begin{array}{l}
                20./\idf{pedestrian\_l},  \\
                20./\idf{pedestrian},     \\
                \bot/\idf{t\_pedestrian}, \\
                \Soft/\idf{strat}
            \end{array}\right]}, ()):
    ([], \, []).
\end{equation*}

\paragraph{Actualize flow operation}

Flow operations are designed to react to event occurrences or signal updates, depending on the
operation, to create the appropriate actualization. Some operations, such as \kwf{persist}, need
to track previous values to detect signal changes and are therefore stateful. Additionally, some
operations, like \kwf{timeout}, are associated with a timer that they reset when necessary.

The actualization of an operation $\sigma$ in the context $\gamma$ is written as
\begin{equation*}
    \sigmaOpSem{\gamma}{s}{\sigma}{t}{s'}{v_\bot^+}{(i_T, t_T)^+}.
\end{equation*}
In this notation, $s$ and $s'$ represent the current and next states, $[v_\bot^+]$ are the
output values, and $[(i_T, t_T)^+]$ are the induced timer resets.

The semantics of the output values depends on the type of flow $\sigma$: if it is a signal, $v_\bot$
represents a signal update ($\bot$ if unchanged and $v$ if changed); if it is an event, $v_\bot$
represents an event occurrence ($\bot$ if absent and $v$ if present).
%
To ensure that signals always have a well-defined value, even before their computations are
triggered, we choose to interpret the first output value as an update to an initial default value
$v_d$ (that is unique and predefined for each type).

\subparagraph{Identifiers}

As described in rule \nameref{rule:semantics:grfrp:opsem:step:op:ident}, an identifier ($x$) returns the value
stored in the update context, independently of the flow kind (signal or event). It then
has an empty initial state and timer queue, as defined by
\Cref{eq:semantics:grfrp:opsem:init:op:ident}.

\begin{align}
    \stateTy{x} & = \{()\}                        & \initstate{x} & = ()                                                         &
    \timerTy{x} & = \List{\emptyset \times \Time} & \inittimer{x} & = \emptylist  \label{eq:semantics:grfrp:opsem:init:op:ident}
\end{align}

\begin{center}
    \begin{minipage}{0.95\textwidth}
        \centering
        \begin{mathpar}
            %%%%% Ident
            \inferOPSEMFlow
            { }
            { \sigmaOpSem{\gamma}{()}{x}{t}{()}{\gamma[x]}{\,} }
            {ID}
            \label{rule:semantics:grfrp:opsem:step:op:ident}
        \end{mathpar}
    \end{minipage}
\end{center}

\subparagraph{Component application}

The operational semantics of \GRSync, detailed in \Cref{sec:semantics:grsync:opsem}, models the
execution of components as state machines: it has an initial state and a step function that
produces, from inputs and the current state, outputs and the next state.

The component application $f(\sigma^+)$ creates a tuple of signals and events, determined by the
specification of the component $f$ in \G. The state machine associated with $f$ by the operational
semantics of \GRSync performs a step only when one of its input signals has changed or one of its
input events has occurred, as depicted in rule \nameref{rule:semantics:grfrp:opsem:step:op:app:value}.
%
The service delivers to the component state machine the current values of all inputs, denoted by
$v_i^+$ and created using the function \smf{into\_val} from \Cref{rule:semantics:grfrp:opsem:aux:%
    intoVal:signal,rule:semantics:grfrp:opsem:aux:intoVal:event}. These inputs include optional
values for events and the most recent values for signals. If a signal has not changed, the service
uses the previous input values stored in $(v_i^\explast)^+$, initialized to $v_d^+$ by \Cref{eq:%
    semantics:grfrp:opsem:init:op:app:state}. The component state machine produces output values
$v_o^+$, which are then filtered to retain only signal changes and event occurrences, represented as
$v_{o\bot}^+$ and created using the function \smf{into\_bot} defined in \Cref{rule:semantics:grfrp:%
    opsem:aux:intoBot:signal,rule:semantics:grfrp:opsem:aux:intoBot:event}. To distinguish signal
changes, the service stores the previous output values in $(v_o^\explast)^+$, also initialized to
$v_d^+$ by \Cref{eq:semantics:grfrp:opsem:init:op:app:state}.
%
When all events and signals return $\bot$, indicating that they are absent or unchanged, the rule
\nameref{rule:semantics:grfrp:opsem:step:op:app:propag} is applied to propagate the flows states.

\begin{align}
    \stateTy{f(\sigma^+)}   & = \stateTy{f} \times \Pi_{x_i} \Val \times \Pi_{x_o} \Val \times \Pi_{\sigma} \stateTy{\sigma} \label{eq:semantics:grfrp:opsem:stateTy:op:app} \\
    \text{with} \;          & \G[f] = \component{f}{x_i^+}{x_o^+}{\dots}{\dots} \nonumber                                                                                    \\
    \initstate{f(\sigma^+)} & = (\initsync{f}, v_d^+, v_d^+, (\initstate{\sigma})^+) \label{eq:semantics:grfrp:opsem:init:op:app:state}                                      \\
    \timerTy{f(\sigma^+)}   & = \List{(\bigcup U^+) \times \Time} \label{eq:semantics:grfrp:opsem:timerTy:op:app}                                                            \\
    \text{with} \;          & \left(\timerTy{\sigma} = \List{U \times \Time}\right)^+ \nonumber                                                                              \\
    \inittimer{f(\sigma^+)} & = \bigLstconcat (\inittimer{\sigma})^+ \label{eq:semantics:grfrp:opsem:init:op:app:timer}
\end{align}

\begin{center}
    \begin{minipage}{0.95\textwidth}
        \centering
        \begin{mathpar}
            %%%%% Component application
            \inferOPSEMFlow
            {
                \G[f] = \component{f}{x_i^+}{x_o^+}{eq^+}{x_m^+ = c^+} \\\\
                \left(\sigmaOpSem{\gamma}{s}{\sigma}{t}{s'}{v_{i\bot}}{(i_\sigma, t_\sigma)^+}\right)^+ \\
                \exists v_{i\bot} \neq \bot \\\\
                (v_i = \intoVal{\tyCtx[x_i]}{v_i^\explast}{v_{i\bot}})^+ \\
                (v_{o\bot} = \intoBot{\tyCtx[x_o]}{v_o^\explast}{v_o})^+ \\\\
                \compOpsem{s_f}{v_i^+}{f}{s'_f}{v_o^+}
            }
            { \sigmaOpSem{\gamma}{(s_f, (v_i^\explast)^+, (v_o^\explast)^+, s^+)}{f(\sigma^+)}{t}
                {(s'_f, v_i^+, v_o^+, s^{\prime+})}{v_{o\bot}^+}{((i_\sigma, t_\sigma)^+)^+} }
            {App{\scriptsize val}}
            \label{rule:semantics:grfrp:opsem:step:op:app:value}

            \inferOPSEMFlow
            { \left(\sigmaOpSem{\gamma}{s}{\sigma}{t}{s'}{\bot}{(i_\sigma, t_\sigma)^+}\right)^+ }
            { \sigmaOpSem{\gamma}{(s_f, (v_i^\explast)^+, (v_o^\explast)^+, s^+)}{f(\sigma^+)}{t}
                {(s_f, (v_i^\explast)^+, (v_o^\explast)^+, s^{\prime+})}{\bot^+}
                {((i_\sigma, t_\sigma)^+)^+} }
            {App{\scriptsize propag}}
            \label{rule:semantics:grfrp:opsem:step:op:app:propag}
        \end{mathpar}
    \end{minipage}
\end{center}

\begin{align}
    \intoVal{ty}{v^\explast}{v_\bot} & = \ttf{match} \, v_\bot \, \{ \, v \rightarrow v, \, \bot \rightarrow v^\explast \, \} \label{rule:semantics:grfrp:opsem:aux:intoVal:signal}  \\
    \intoVal{\mayB{ty}}{\_}{v_\bot}  & = \ttf{match} \, v_\bot \, \{ \, v \rightarrow \some{v}, \, \bot \rightarrow \none \, \} \label{rule:semantics:grfrp:opsem:aux:intoVal:event} \\
    \intoBot{ty}{v^\explast}{v}      & = \text{if} \, v \neq v^\explast \, \text{then} \, v \, \text{else} \, \bot \label{rule:semantics:grfrp:opsem:aux:intoBot:signal}             \\
    \intoBot{\mayB{ty}}{\_}{v}       & = \ttf{match} \, v \, \{ \, \some{v'} \rightarrow v' \, \none \rightarrow \bot \, \} \label{rule:semantics:grfrp:opsem:aux:intoBot:event}
\end{align}

\subparagraph{Sample}

The operation \sample{\sigma}{T} creates an event that emits the most recent value of the flow
$\sigma$ every $T$ seconds, provided that $\sigma$ has occurred. This gives control over the
timing of the event, dictating its occurrence according to the specified period. In the semantics,
the occurrence of the timer is verified using the update context: if $\gamma[i_T] = ()$, then the
timer associated with \sample{\sigma}{T} has occurred; if $\gamma[i_T] = \bot$, then its timer is
absent.
%
To achieve this, the operation needs to store the previous value of $\sigma$ in its state,
represented as a bottom value $v_\bot \in \ValBot$ initialized to $\bot$ in
\Cref{eq:semantics:grfrp:opsem:init:op:sample:state}.
%
If $\sigma$ has not occurred since the last period, $v_\bot$ is set to $\bot$. When $\sigma$ has a
value, the rule \nameref{rule:semantics:grfrp:opsem:step:op:sample:value} is applied to update the memory $v_\bot$ with the
new value $v$. If $\sigma$ returns $\bot$, indicating that the event is absent, the rule
\nameref{rule:semantics:grfrp:opsem:step:op:sample:propag} propagates the new state $s'$ of $\sigma$. If the
propagation is triggered by a timer $i_T$ associated with the \kwf{sample} operation, the rule
\nameref{rule:semantics:grfrp:opsem:step:op:sample:timer} states that the operation evaluates to the stored
value $v$, the memory is then reset to $\bot$, and the corresponding timer is reset.
%
In order to initiate this timer, it is added to the initial timer queue by
\Cref{eq:semantics:grfrp:opsem:init:op:sample:timer}.

\begin{align}
    \stateTy{\sample{\sigma}{T}}   & = \ValBot \times \stateTy{\sigma} \label{eq:semantics:grfrp:opsem:stateTy:op:sample}      \\
    \initstate{\sample{\sigma}{T}} & = (\bot, \initstate{\sigma}) \label{eq:semantics:grfrp:opsem:init:op:sample:state}        \\
    \timerTy{\sample{\sigma}{T}}   & = \List{(\{i_T\} \cup U) \times \Time} \label{eq:semantics:grfrp:opsem:timerTy:op:sample} \\
    \text{with} \;                 & \timerTy{\sigma} = \List{U \times \Time} \nonumber                                        \\
    \inittimer{\sample{\sigma}{T}} & = (i_T, T) :: \inittimer{\sigma} \label{eq:semantics:grfrp:opsem:init:op:sample:timer}
\end{align}

\begin{center}
    \begin{minipage}{0.95\textwidth}
        \centering
        \begin{mathpar}
            %%%%% Sample
            \inferOPSEMFlow
            {\sigmaOpSem{\gamma}{s}{\sigma}{t}{s'}{v}{(i_\sigma, t_\sigma)^+}}
            {\sigmaOpSem{\gamma}{(v_\bot, s)}{\sample{\sigma}{T}}{t}{(v, s')}{\bot}{(i_\sigma, t_\sigma)^+}}
            {Samp{\scriptsize val}}
            \label{rule:semantics:grfrp:opsem:step:op:sample:value}

            \inferOPSEMFlow
            {
                \gamma[i_T] = \bot \\
                \sigmaOpSem{\gamma}{s}{\sigma}{t}{s'}{\bot}{(i_\sigma, t_\sigma)^+}
            }
            {\sigmaOpSem{\gamma}{(v_\bot, s)}{\sample{\sigma}{T}}{t}{(v_\bot, s')}{\bot}{(i_\sigma, t_\sigma)^+}}
            {Samp{\scriptsize propag}}
            \label{rule:semantics:grfrp:opsem:step:op:sample:propag}

            \inferOPSEMFlow
            {
                \gamma[i_T] = () \\
                \sigmaOpSem{\gamma}{s}{\sigma}{t}{s'}{\bot}{(i_\sigma, t_\sigma)^+}
            }
            {\sigmaOpSem{\gamma}{(v_\bot, s)}{\sample{\sigma}{T}}{t}{(\bot, s')}{v_\bot}{(i_T, t + T), (i_\sigma, t_\sigma)^+}}
            {Samp{\scriptsize timer}}
            \label{rule:semantics:grfrp:opsem:step:op:sample:timer}
        \end{mathpar}
    \end{minipage}
\end{center}

\subparagraph{Scan}

The flow operation \scan{\sigma}{T} generates a signal that only includes the value updates from
$\sigma$ at each period $T$. It stores the most recent value update of $\sigma$ named
$v_\bot \in \ValBot$, as well as the last value of \scan{\sigma}{T} named $v^\explast \in \Val$,
respectively initialized to $\bot$ and a default value $v_d$ in
\Cref{eq:semantics:grfrp:opsem:init:op:scan:state}.
%
The different possible steps represent all the combinations of three parameters: the presence or
absence of the associated timer $i_T$, the possible update of the input signal $\sigma$, and
if the update of $\sigma$ is different from $v^\explast$ (\ie if it represents an update for
\scan{\sigma}{T}).
%
When the timer $i_T$ associated with the \kwf{scan} operation is absent and $\sigma$ has a value
different from $v^\explast$, the rule \nameref{rule:semantics:grfrp:opsem:step:op:scan:value} is
applied to update the memory $v_\bot$ with the new value. If, however, the value of $\sigma$ is
found to be equal to $v^\explast$, the rule \nameref{rule:semantics:grfrp:opsem:step:op:scan:drop}
resets $v_\bot$ to $\bot$, indicating that no update has been made to the signal \scan{\sigma}{T}.
If $i_T$ is absent and $\sigma$ returns $\bot$, indicating that the signal is unchanged, the rule
\nameref{rule:semantics:grfrp:opsem:step:op:scan:propag} propagates the new state $s'$ of $\sigma$.
If the propagation is triggered by $i_T$ and caused the update of $\sigma$ to a value $v$ different
from $v^\explast$, the rule \nameref{rule:semantics:grfrp:opsem:step:op:scan:valuetimer} ensures
that \scan{\sigma}{T} is updated to $v$. In this case, the memory is reset to $\bot$, the last value
is updated to $v$, and the corresponding timer is reset. If the update of $\sigma$ is equal to
$v^\explast$ then it is an update for \scan{\sigma}{T}, which returns $\bot$, resets its memory to
$\bot$ and the corresponding timer, as defined by the rule
\nameref{rule:semantics:grfrp:opsem:step:op:scan:droptimer}. If $i_T$ is present, $\sigma$ returns
$\bot$, and the stored value is $v$, the rule
\nameref{rule:semantics:grfrp:opsem:step:op:scan:timer} specifies that the operation evaluates to
$v$, the memory is reset to $\bot$, the last value is updated to $v$, and the corresponding timer is
reset. If, instead, the stored value is $\bot$, the rule
\nameref{rule:semantics:grfrp:opsem:step:op:scan:bot} indicates that the operation evaluates to
$\bot$. Here, the memory is reset to $\bot$, the last value remains unchanged, and the corresponding
timer is reset.
%
Like \kwf{sample}, the timer is added to the initial timer queue by \Cref{eq:semantics:grfrp:opsem:init:op:scan:timer}.

\begin{align}
    \stateTy{\scan{\sigma}{T}}   & = \ValBot \times \Val \times \stateTy{\sigma} \label{eq:semantics:grfrp:opsem:stateTy:op:scan} \\
    \initstate{\scan{\sigma}{T}} & = (\bot, v_d, \initstate{\sigma}) \label{eq:semantics:grfrp:opsem:init:op:scan:state}          \\
    \timerTy{\scan{\sigma}{T}}   & = \List{(\{i_T\} \cup U) \times \Time} \label{eq:semantics:grfrp:opsem:timerTy:op:scan}        \\
    \text{with} \;               & \timerTy{\sigma} = \List{U \times \Time} \nonumber                                             \\
    \inittimer{\scan{\sigma}{T}} & = (i_T, T) :: \inittimer{\sigma} \label{eq:semantics:grfrp:opsem:init:op:scan:timer}
\end{align}

\begin{center}
    \begin{minipage}{0.95\textwidth}
        \centering
        \begin{mathpar}
            %%%%% Scan
            \inferOPSEMFlow
            {
                \gamma[i_T] = \bot \\
                \sigmaOpSem{\gamma}{s}{\sigma}{t}{s'}{v}{(i_\sigma, t_\sigma)^+} \\
                v \neq v^\explast
            }
            {\sigmaOpSem{\gamma}{(v_\bot, v^\explast, s)}{\scan{\sigma}{T}}{t}{(v, v^\explast, s')}{\bot}{(i_\sigma, t_\sigma)^+}}
            {Scan{\scriptsize val}}
            \label{rule:semantics:grfrp:opsem:step:op:scan:value}

            \inferOPSEMFlow
            {
                \gamma[i_T] = \bot \\
                \sigmaOpSem{\gamma}{s}{\sigma}{t}{s'}{v}{(i_\sigma, t_\sigma)^+} \\
                v = v^\explast
            }
            {\sigmaOpSem{\gamma}{(v_\bot, v^\explast, s)}{\scan{\sigma}{T}}{t}{(\bot, v^\explast, s')}{\bot}{(i_\sigma, t_\sigma)^+}}
            {Scan{\scriptsize drop}}
            \label{rule:semantics:grfrp:opsem:step:op:scan:drop}

            \inferOPSEMFlow
            {
                \gamma[i_T] = \bot \\
                \sigmaOpSem{\gamma}{s}{\sigma}{t}{s'}{\bot}{(i_\sigma, t_\sigma)^+}
            }
            {\sigmaOpSem{\gamma}{(v_\bot, v^\explast, s)}{\scan{\sigma}{T}}{t}{(v_\bot, v^\explast, s')}{\bot}{(i_\sigma, t_\sigma)^+}}
            {Scan{\scriptsize propag}}
            \label{rule:semantics:grfrp:opsem:step:op:scan:propag}

            \inferOPSEMFlow
            {
                \gamma[i_T] = () \\
                \sigmaOpSem{\gamma}{s}{\sigma}{t}{s'}{v}{(i_\sigma, t_\sigma)^+} \\
                v \neq v^\explast
            }
            {\sigmaOpSem{\gamma}{(v_\bot, v^\explast, s)}{\scan{\sigma}{T}}{t}{(\bot, v, s')}{v}{(i_T, t + T), (i_\sigma, t_\sigma)^+}}
            {Scan{\scriptsize valTimer}}
            \label{rule:semantics:grfrp:opsem:step:op:scan:valuetimer}

            \inferOPSEMFlow
            {
                \gamma[i_T] = () \\
                \sigmaOpSem{\gamma}{s}{\sigma}{t}{s'}{v}{(i_\sigma, t_\sigma)^+} \\
                v = v^\explast
            }
            {\sigmaOpSem{\gamma}{(v_\bot, v^\explast, s)}{\scan{\sigma}{T}}{t}{(\bot, v^\explast, s')}{\bot}{(i_T, t + T), (i_\sigma, t_\sigma)^+}}
            {Scan{\scriptsize dropTimer}}
            \label{rule:semantics:grfrp:opsem:step:op:scan:droptimer}

            \inferOPSEMFlow
            {
                \gamma[i_T] = () \\
                \sigmaOpSem{\gamma}{s}{\sigma}{t}{s'}{\bot}{(i_\sigma, t_\sigma)^+}
            }
            {\sigmaOpSem{\gamma}{(v, v^\explast, s)}{\scan{\sigma}{T}}{t}{(\bot, v, s')}{v}{(i_T, t + T), (i_\sigma, t_\sigma)^+}}
            {Scan{\scriptsize timer}}
            \label{rule:semantics:grfrp:opsem:step:op:scan:timer}

            \inferOPSEMFlow
            {
                \gamma[i_T] = () \\
                \sigmaOpSem{\gamma}{s}{\sigma}{t}{s'}{\bot}{(i_\sigma, t_\sigma)^+}
            }
            {\sigmaOpSem{\gamma}{(\bot, v^\explast, s)}{\scan{\sigma}{T}}{t}{(\bot, v^\explast, s')}{\bot}{(i_T, t + T), (i_\sigma, t_\sigma)^+}}
            {Scan{\scriptsize bot}}
            \label{rule:semantics:grfrp:opsem:step:op:scan:bot}
        \end{mathpar}
    \end{minipage}
\end{center}

\subparagraph{Timeout}

The flow operation \timeout{\sigma}{T} creates an event with unit values, triggered when the event
$\sigma$ has not emitted a value for more than $T$ seconds. When $\sigma$ has a value, the rule
\nameref{rule:semantics:grfrp:opsem:step:op:timeout:reset} is applied to reset the corresponding timer.
If $\sigma$ returns $\bot$, indicating the event is absent, the rule
\nameref{rule:semantics:grfrp:opsem:step:op:timeout:propag} propagates the new state $s'$ of $\sigma$.
If the propagation is triggered by a timer $i_T$ associated with the \kwf{timeout} operation, the
rule \nameref{rule:semantics:grfrp:opsem:step:op:timeout:timer} specifies that the operation evaluates
to $()$, and the corresponding timer is reset. Like \kwf{sample} and \kwf{scan}, the timer is
added to the initial timer queue by \Cref{eq:semantics:grfrp:opsem:init:op:timeout:timer}.

\begin{align}
    \stateTy{\timeout{\sigma}{T}}   & = \stateTy{\sigma} \label{eq:semantics:grfrp:opsem:stateTy:op:timeout}                     \\
    \initstate{\timeout{\sigma}{T}} & = \initstate{\sigma} \label{eq:semantics:grfrp:opsem:init:op:timeout:state}                \\
    \timerTy{\timeout{\sigma}{T}}   & = \List{(\{i_T\} \cup U) \times \Time} \label{eq:semantics:grfrp:opsem:timerTy:op:timeout} \\
    \text{with} \;                  & \timerTy{\sigma} = \List{U \times \Time} \nonumber                                         \\
    \inittimer{\timeout{\sigma}{T}} & = (i_T, T) :: \inittimer{\sigma} \label{eq:semantics:grfrp:opsem:init:op:timeout:timer}
\end{align}

\begin{center}
    \begin{minipage}{0.95\textwidth}
        \centering
        \begin{mathpar}
            %%%%% Timeout
            \inferOPSEMFlow
            { \sigmaOpSem{\gamma}{s}{\sigma}{t}{s'}{v}{(i_\sigma, t_\sigma)^+} }
            { \sigmaOpSem{\gamma}{s}{\timeout{\sigma}{T}}{t}{s'}{\bot}{(i_T, t + T), (i_\sigma, t_\sigma)^+} }
            {TOut{\scriptsize reset}}
            \label{rule:semantics:grfrp:opsem:step:op:timeout:reset}

            \inferOPSEMFlow
            {
                \gamma[i_T] = \bot \\
                \sigmaOpSem{\gamma}{s}{\sigma}{t}{s'}{\bot}{(i_\sigma, t_\sigma)^+}
            }
            { \sigmaOpSem{\gamma}{s}{\timeout{\sigma}{T}}{t}{s'}{\bot}{(i_\sigma, t_\sigma)^+} }
            {TOut{\scriptsize propag}}
            \label{rule:semantics:grfrp:opsem:step:op:timeout:propag}

            \inferOPSEMFlow
            {
                \gamma[i_T] = () \\
                \sigmaOpSem{\gamma}{s}{\sigma}{t}{s'}{\bot}{(i_\sigma, t_\sigma)^+}
            }
            { \sigmaOpSem{\gamma}{s}{\timeout{\sigma}{T}}{t}{s'}{()}{(i_T, t + T), (i_\sigma, t_\sigma)^+} }
            {TOut{\scriptsize timer}}
            \label{rule:semantics:grfrp:opsem:step:op:timeout:timer}
        \end{mathpar}
    \end{minipage}
\end{center}

\subparagraph{Throttle}

The flow operation \throttle{\sigma}{\Delta} creates a signal that retains only the values from
$\sigma$ whose difference from the previous value exceeds $\Delta$. To propagate updates for
\throttle{\sigma}{\Delta}, its last value, denoted by $v^\explast \in \Val$, must be stored and
initialized to a default value $v_d$, as defined in
\Cref{eq:semantics:grfrp:opsem:stateTy:op:throttle,eq:semantics:grfrp:opsem:init:op:throttle:state}.
When $\sigma$ returns a value such that the difference from $v^\explast$ exceeds $\Delta$, the rule
\nameref{rule:semantics:grfrp:opsem:step:op:throttle:value} is applied, returning this new value as
output and storing it as the updated $v^\explast$. Conversely, if the difference from $v^\explast$
is smaller than $\Delta$, the rule \nameref{rule:semantics:grfrp:opsem:step:op:throttle:drop} is
applied, dropping the value, returning $\bot$, and leaving $v^\explast$ unchanged. Finally, if
$\sigma$ returns $\bot$, indicating that the signal remains unchanged, the rule
\nameref{rule:semantics:grfrp:opsem:step:op:throttle:propag} is applied to propagate only the
updated state $s'$ of $\sigma$.

\begin{align}
    \stateTy{\throttle{\sigma}{\Delta}}   & = \Val \times \stateTy{\sigma} \label{eq:semantics:grfrp:opsem:stateTy:op:throttle} \\
    \initstate{\throttle{\sigma}{\Delta}} & = (v_d, \initstate{\sigma}) \label{eq:semantics:grfrp:opsem:init:op:throttle:state} \\
    \timerTy{\throttle{\sigma}{\Delta}}   & = \timerTy{\sigma} \label{eq:semantics:grfrp:opsem:timerTy:op:throttle}             \\
    \inittimer{\throttle{\sigma}{\Delta}} & = \inittimer{\sigma} \label{eq:semantics:grfrp:opsem:init:op:throttle:timer}
\end{align}

\begin{center}
    \begin{minipage}{0.95\textwidth}
        \centering
        \begin{mathpar}
            %%%%% Throttle
            \inferOPSEMFlow
            {
                \sigmaOpSem{\gamma}{s}{\sigma}{t}{s'}{v}{(i_\sigma, t_\sigma)^+} \\
                |v - v^\explast| \ge \Delta
            }
            { \sigmaOpSem{\gamma}{(v^\explast, s)}{\throttle{\sigma}{\Delta}}{t}{(v, s')}{v}{(i_\sigma, t_\sigma)^+} }
            {Thrt{\scriptsize val}}
            \label{rule:semantics:grfrp:opsem:step:op:throttle:value}

            \inferOPSEMFlow
            {
                \sigmaOpSem{\gamma}{s}{\sigma}{t}{s'}{v}{(i_\sigma, t_\sigma)^+} \\
                |v - v^\explast| < \Delta
            }
            { \sigmaOpSem{\gamma}{(v^\explast, s)}{\throttle{\sigma}{\Delta}}{t}{(v^\explast, s')}{\bot}{(i_\sigma, t_\sigma)^+} }
            {Thrt{\scriptsize drop}}
            \label{rule:semantics:grfrp:opsem:step:op:throttle:drop}

            \inferOPSEMFlow
            { \sigmaOpSem{\gamma}{s}{\sigma}{t}{s'}{\bot}{(i_\sigma, t_\sigma)^+} }
            { \sigmaOpSem{\gamma}{(v^\explast, s)}{\throttle{\sigma}{\Delta}}{t}{(v^\explast, s')}{\bot}{(i_\sigma, t_\sigma)^+} }
            {Thrt{\scriptsize propag}}
            \label{rule:semantics:grfrp:opsem:step:op:throttle:propag}
        \end{mathpar}
    \end{minipage}
\end{center}

\subparagraph{Onchange}

The flow operation \onchange{\sigma} derives an event from a signal. The values of the output event
are identical to those of the input signal, as specified by the rule
\nameref{rule:semantics:grfrp:opsem:step:op:onchange}. However, the interpretation and behavior of
these values differ: for the output event, they represent event occurrences, whereas for the input
signal, they represent signal updates. The initial state and timer queue of \onchange{\sigma} are
inherited from $\sigma$, as defined in \Crefrange{eq:semantics:grfrp:opsem:init:op:onchange:state}
{eq:semantics:grfrp:opsem:init:op:onchange:timer}.

\begin{align}
    \stateTy{\onchange{\sigma}}   & = \stateTy{\sigma} \label{eq:semantics:grfrp:opsem:stateTy:op:onchange}      \\
    \initstate{\onchange{\sigma}} & = \initstate{\sigma} \label{eq:semantics:grfrp:opsem:init:op:onchange:state} \\
    \timerTy{\onchange{\sigma}}   & = \timerTy{\sigma} \label{eq:semantics:grfrp:opsem:timerTy:op:onchange}      \\
    \inittimer{\onchange{\sigma}} & = \inittimer{\sigma} \label{eq:semantics:grfrp:opsem:init:op:onchange:timer}
\end{align}

\begin{center}
    \begin{minipage}{0.95\textwidth}
        \centering
        \begin{mathpar}
            %%%%% Onchange
            \inferOPSEMFlow
            { \sigmaOpSem{\gamma}{s}{\sigma}{t}{s'}{v_\bot}{(i_\sigma, t_\sigma)^+} }
            { \sigmaOpSem{\gamma}{s}{\onchange{\sigma}}{t}{s'}{v_\bot}{(i_\sigma, t_\sigma)^+} }
            {OChg}
            \label{rule:semantics:grfrp:opsem:step:op:onchange}
        \end{mathpar}
    \end{minipage}
\end{center}

\subparagraph{Persist}

The flow operation \persist{\sigma} mirrors the \kwf{on\_change} operation, but in the opposite
direction: it derives a signal from an event. To do so, it filters out redundant values, as signal
output values must represent meaningful updates. The last value of the created signal, denoted by
$v^\explast \in \Val$, is stored and initialized to a default value $v_d$, as defined in
\Cref{eq:semantics:grfrp:opsem:init:op:persist:state}. This filtering behavior is governed by the
rule \nameref{rule:semantics:grfrp:opsem:step:op:persist:drop}, which applies when the input event
$\sigma$ emits a value equal to $v^\explast$. In such cases, the value is dropped, and
\persist{\sigma} returns $\bot$, ensuring that only meaningful updates are retained. In all other
cases, the values of $\sigma$ are directly passed to \persist{\sigma}, as described by the rules
\nameref{rule:semantics:grfrp:opsem:step:op:persist:value} and
\nameref{rule:semantics:grfrp:opsem:step:op:persist:propag}.

\begin{align}
    \stateTy{\persist{\sigma}}   & = \Val \times \stateTy{\sigma} \label{eq:semantics:grfrp:opsem:stateTy:op:persist} \\
    \initstate{\persist{\sigma}} & = (v_d, \initstate{\sigma}) \label{eq:semantics:grfrp:opsem:init:op:persist:state} \\
    \timerTy{\persist{\sigma}}   & = \timerTy{\sigma} \label{eq:semantics:grfrp:opsem:timerTy:op:persist}             \\
    \inittimer{\persist{\sigma}} & = \inittimer{\sigma} \label{eq:semantics:grfrp:opsem:init:op:persist:timer}
\end{align}

\begin{center}
    \begin{minipage}{0.95\textwidth}
        \centering
        \begin{mathpar}
            %%%%% Persist
            \inferOPSEMFlow
            {
                \sigmaOpSem{\gamma}{s}{\sigma}{t}{s'}{v}{(i_\sigma, t_\sigma)^+} \\
                v = v^\explast
            }
            { \sigmaOpSem{\gamma}{(v^\explast, s)}{\persist{\sigma}}{t}{(v^\explast, s')}{\bot}{(i_\sigma, t_\sigma)^+} }
            {Prst{\scriptsize drop}}
            \label{rule:semantics:grfrp:opsem:step:op:persist:drop}

            \inferOPSEMFlow
            {
                \sigmaOpSem{\gamma}{s}{\sigma}{t}{s'}{v}{(i_\sigma, t_\sigma)^+} \\
                v \neq v^\explast
            }
            { \sigmaOpSem{\gamma}{(v^\explast, s)}{\persist{\sigma}}{t}{(v, s')}{v}{(i_\sigma, t_\sigma)^+} }
            {Prst{\scriptsize val}}
            \label{rule:semantics:grfrp:opsem:step:op:persist:value}

            \inferOPSEMFlow
            { \sigmaOpSem{\gamma}{s}{\sigma}{t}{s'}{\bot}{(i_\sigma, t_\sigma)^+} }
            { \sigmaOpSem{\gamma}{(v^\explast, s)}{\persist{\sigma}}{t}{(v^\explast, s')}{\bot}{(i_\sigma, t_\sigma)^+} }
            {Prst{\scriptsize propag}}
            \label{rule:semantics:grfrp:opsem:step:op:persist:propag}
        \end{mathpar}
    \end{minipage}
\end{center}

\subparagraph{Merge}

The flow operation $\mergeop{\sigma_1}{\sigma_2}$ combines the values from the events $\sigma_1$ and
$\sigma_2$ into a new event. In cases where both events occur simultaneously, as handled by rule
\nameref{rule:semantics:grfrp:opsem:step:op:merge:1}, the resulting event prioritizes the values from
$\sigma_1$. If only one event occurs, rules \nameref{rule:semantics:grfrp:opsem:step:op:merge:1} or
\nameref{rule:semantics:grfrp:opsem:step:op:merge:2} are applied, depending on which event occurred, and
the value of this event is returned. When both events return $\bot$, indicating their absence, rule
\nameref{rule:semantics:grfrp:opsem:step:op:merge:propag} propagates the updated states $s_1'$ and $s_2'$
of $\sigma_1$ and $\sigma_2$, respectively. Its initial state and timer queue correspond to those of
$\sigma_1$ and $\sigma_2$, as defined in \Cref{eq:semantics:grfrp:opsem:init:op:merge:timer}.

\begin{align}
    \stateTy{\mergeop{\sigma_1}{\sigma_2}}   & = \stateTy{\sigma_1} \times \stateTy{\sigma_2} \label{eq:semantics:grfrp:opsem:stateTy:op:merge}                     \\
    \initstate{\mergeop{\sigma_1}{\sigma_2}} & = (\initstate{\sigma_1}, \initstate{\sigma_2}) \label{eq:semantics:grfrp:opsem:init:op:merge:state}                  \\
    \timerTy{\mergeop{\sigma_1}{\sigma_2}}   & = \List{(U_1 \cup U_2) \times \Time} \label{eq:semantics:grfrp:opsem:timerTy:op:merge}                               \\
    \text{with} \;                           & \timerTy{\sigma_1} = \List{U_1 \times \Time} \; \text{and} \; \timerTy{\sigma_2} = \List{U_2 \times \Time} \nonumber \\
    \inittimer{\mergeop{\sigma_1}{\sigma_2}} & = \inittimer{\sigma_1} \lstconcat \inittimer{\sigma_2} \label{eq:semantics:grfrp:opsem:init:op:merge:timer}
\end{align}

\begin{center}
    \begin{minipage}{0.95\textwidth}
        \centering
        \begin{mathpar}
            %%%% Merge
            \inferOPSEMFlow
            {
                \sigmaOpSem{\gamma}{s_1}{\sigma_1}{t}{s'_1}{v_1}{(i_{\sigma_1}, t_{\sigma_1})^+} \\
                \sigmaOpSem{\gamma}{s_2}{\sigma_2}{t}{s'_2}{v_{2\bot}}{(i_{\sigma_2}, t_{\sigma_2})^+}
            }
            { \sigmaOpSem{\gamma}{(s_1, s_2)}{\mergeop{\sigma_1}{\sigma_2}}{t}{(s'_1, s'_2)}{v_1}{(i_{\sigma_1}, t_{\sigma_1})^+, (i_{\sigma_2}, t_{\sigma_2})^+} }
            {Mrg{\scriptsize val1}}
            \label{rule:semantics:grfrp:opsem:step:op:merge:1}

            \inferOPSEMFlow
            {
                \sigmaOpSem{\gamma}{s_1}{\sigma_1}{t}{s'_1}{\bot}{(i_{\sigma_1}, t_{\sigma_1})^+} \\
                \sigmaOpSem{\gamma}{s_2}{\sigma_2}{t}{s'_2}{v_2}{(i_{\sigma_2}, t_{\sigma_2})^+}
            }
            { \sigmaOpSem{\gamma}{(s_1, s_2)}{\mergeop{\sigma_1}{\sigma_2}}{t}{(s'_1, s'_2)}{v_2}{(i_{\sigma_1}, t_{\sigma_1})^+, (i_{\sigma_2}, t_{\sigma_2})^+} }
            {Mrg{\scriptsize val2}}
            \label{rule:semantics:grfrp:opsem:step:op:merge:2}

            \inferOPSEMFlow
            {
                \sigmaOpSem{\gamma}{s_1}{\sigma_1}{t}{s'_1}{\bot}{(i_{\sigma_1}, t_{\sigma_1})^+} \\
                \sigmaOpSem{\gamma}{s_2}{\sigma_2}{t}{s'_2}{\bot}{(i_{\sigma_2}, t_{\sigma_2})^+}
            }
            { \sigmaOpSem{\gamma}{(s_1, s_2)}{\mergeop{\sigma_1}{\sigma_2}}{t}{(s'_1, s'_2)}{\bot}{(i_{\sigma_1}, t_{\sigma_1})^+, (i_{\sigma_2}, t_{\sigma_2})^+} }
            {Mrg{\scriptsize propag}}
            \label{rule:semantics:grfrp:opsem:step:op:merge:propag}
        \end{mathpar}
    \end{minipage}
\end{center}

\setcounter{subparagraph}{0}
\subparagraph{Time}

The flow operation \timeop represents time as a signal. During its execution, it returns only the
timestamp of the current propagation, as specified by the rule
\nameref{rule:semantics:grfrp:opsem:step:op:time}. It then has an empty initial state and timer queue, as
defined by \Cref{eq:semantics:grfrp:opsem:init:op:time:state,eq:semantics:grfrp:opsem:init:op:time:timer}.

\begin{align}
    \stateTy{\timeop} & = \{()\}                        & \initstate{\timeop} & = () \label{eq:semantics:grfrp:opsem:init:op:time:state}         \\
    \timerTy{\timeop} & = \List{\emptyset \times \Time} & \inittimer{\timeop} & = \emptylist \label{eq:semantics:grfrp:opsem:init:op:time:timer}
\end{align}

\begin{center}
    \begin{minipage}{0.95\textwidth}
        \centering
        \begin{mathpar}
            %%%%% Time
            \inferOPSEMFlow
            { }
            {\sigmaOpSem{\gamma}{()}{\timeop}{t}{()}{t}{\,}}
            {Time}
            \label{rule:semantics:grfrp:opsem:step:op:time}
        \end{mathpar}
    \end{minipage}
\end{center}

\paragraph{Summary}
This section introduced the operational semantics for \GRFRP services, formalizing their push-based,
event-driven execution model. The semantics models how a service reacts to external messages%
---either a new value or a timer expiration---by initiating a deterministic propagation cascade.
The propagation context tracks signal updates and event occurrences within a single reaction step.
We defined a hierarchy of transition rules: a top-level rule for service message handling, a
recursive rule for propagating updates through statements, and a set of rules for evaluating
individual flow operations. The semantics detailed the behavior of stateful flow operations,
including the interface to synchronous \GRSync components and the management of timers for
time-based operators like \kwf{timeout}, \kwf{sample}, and \kwf{scan}. Together, these rules provide
a rigorous foundation for the asynchronous, message-passing behavior of \GRFRP, ensuring that
computations are performed efficiently only when necessary.
